\section*{第 \ref{ch:g12:reaction-rates} 章——反应速率}

\begin{solution}{exo:g12:reaction-rates:1}
温度；接触面积；浓度；温度。
\end{solution}

\begin{solution}{exo:g12:reaction-rates:2}
\qty{13.9}{min}（蓝色）和 \qty{6.9}{min}（红色）。经过两个\omterm{def:g12:reaction-rates:half-life}{半衰期}，
$10.0 / 4 = \qty{2.5}{mmol/L}$。
\end{solution}

\begin{solution}{exo:g12:reaction-rates:3}
斜率 $(5.0 - 1.0)/10 = \qty{0.40}{mmol.L^{-1}.min^{-1}}$，即
\qty{5}{min} 时的\omterm{def:g12:reaction-rates:rate}{生成速率}。
\end{solution}

\begin{solution}{exo:g12:reaction-rates:4}
降温并\omterm{def:g10:concentration-and-dilution:dilution}{稀释}后，反应变得极慢，以致\omterm{def:g11:titration:titration}{滴定}过程中
组成不再发生变化。所测得的组成，
即为倒入时刻反应被终止时\omterm{def:g6:pure-substances-and-mixtures:mixture}{混合物}的组成。
\end{solution}

\begin{solution}{exo:g12:reaction-rates:5}
\qty{30}{min} 为 3 个\omterm{def:g12:reaction-rates:half-life}{半衰期}，剩余 $1/8$。\qty{1}{h} 为 6 个\omterm{def:g12:reaction-rates:half-life}{半衰期}：
$1/64$，约 \qty{1.6}{\percent}。
\end{solution}

\begin{solution}{exo:g12:reaction-rates:6}
$\ln[\ce{A}]$：2.079、1.733、1.386、1.040、0.693。
每隔 \qty{5}{min} 下降 0.346，图像为直线，属\omterm{def:g12:reaction-rates:first-order}{一级反应}。$k = 0.346/5 =
\qty{0.0693}{min^{-1}}$，且 $t_{1/2} = \ln 2 / k = \qty{10.0}{min}$（
由数据直接可见：每隔 \qty{10}{min} 依次为 8.00、4.00、2.00）。
\end{solution}

\begin{solution}{exo:g12:reaction-rates:7}
$k = 0.693 / 25 = \qty{0.0277}{s^{-1}}$。
\end{solution}

\begin{solution}{exo:g12:reaction-rates:8}
$0.100\,e^{-0.20} = \qty{0.0819}{mol/L}$；$0.100\,e^{-0.60} =
\qty{0.0549}{mol/L}$；$0.100\,e^{-2.0} = \qty{0.0135}{mol/L}$。
\end{solution}

\begin{solution}{exo:g12:reaction-rates:9}
在 \qty{13.9}{min} 时，$[\ce{A}] = \qty{5.0}{mmol/L}$；切线斜率接近
$-0.25$，且 $-k[\ce{A}] = -0.050 \times 5.0 =
\qty{-0.25}{mmol.L^{-1}.min^{-1}}$：恰为初始值的一半。
\end{solution}

\begin{solution}{exo:g12:reaction-rates:10}
碘是唯一有色的物种。由比尔–朗伯定律，
$A = \varepsilon \ell [\ce{I2}]$。\omterm{def:g12:reaction-rates:rate}{生成速率}即为
曲线 $[\ce{I2}](t)$ 的切线斜率，即
$A(t)$ 的切线斜率除以 $\varepsilon \ell$。
\end{solution}

\begin{solution}{exo:g12:reaction-rates:11}
第二次实验的初始速率 $k[\ce{A}]_0$ 为第一次的两倍；
两次的\omterm{def:g12:reaction-rates:half-life}{半衰期}相等，即 $\ln 2 / k$，与 $[\ce{A}]_0$ 无关。
\end{solution}

\begin{solution}{exo:g12:reaction-rates:12}
$t = \ln(1/f)/k$。以\omterm{def:g12:reaction-rates:half-life}{半衰期}计为 $\ln(1/f)/\ln 2$。对于 \qty{1}{\percent}，
$\ln 100 / \ln 2 = 6.6$；对于 \qty{0.1}{\percent}，$\ln 1000 / \ln 2 =
10.0$。
\end{solution}

\begin{solution}{exo:g12:reaction-rates:13}
\qty{90}{\percent} 已消耗，$f = 0.10$：$t = \ln 10 / k$，即
\qty{46}{min}（\qty{20}{\celsius} 时），\qty{23}{min}（\qty{30}{\celsius} 时）。
\omterm{def:g12:reaction-rates:first-order}{速率常数}每升高 \qty{10}{\celsius} 约加倍，本题符合这一规律。
\end{solution}

\begin{solution}{exo:g12:reaction-rates:14}
温度降至接近 \qty{0}{\celsius}，且浓度
均降为原来的十分之一。速率正比于两个浓度之积，
故仅\omterm{def:g10:concentration-and-dilution:dilution}{稀释}就使速率除以 $10 \times 10 = 100$。
\end{solution}

\begin{solution}{exo:g12:reaction-rates:15}
$v(t) = -\dfrac{d}{dt}\big([\ce{A}]_0 e^{-kt}\big) = k[\ce{A}]_0 e^{-kt}$。
初始速率为 $k[\ce{A}]_0$；$e^{-kt} = 1/2$ 时速率减半，即
经过一个\omterm{def:g12:reaction-rates:half-life}{半衰期}。
\end{solution}

\begin{solution}{pb:g12:reaction-rates:1}
\textbf{1.} 由比尔–朗伯定律，染料是此波长下唯一吸光的物种，
溶液必须足够稀，
（$A$ 约低于 1）。

\textbf{2.} 它不吸收可见光，对 $A$ 没有贡献。

\textbf{3.} 为使其浓度几乎不变，此时速率
仅取决于染料。

\textbf{4.} $A = 0$（空白对照）。

\textbf{5.} $A$ 在 \qty{6}{min} 内由 0.800 降至 0.402，约为
\qty{6}{min}。

\textbf{6.} 由 \qty{6}{min} 时的 0.402 降至 \qty{12}{min} 时的 0.199，
又经 \qty{6}{min} 减半。

\textbf{7.} $-0.223$、$-0.448$、$-0.691$、$-0.911$、$-1.155$、$-1.366$、
$-1.614$、$-1.945$、$-2.551$、$-3.079$。

\textbf{8.} 各点位于一条直线上，故反应
对染料为\omterm{def:g12:reaction-rates:first-order}{一级反应}。

\textbf{9.} $(-2.551 + 0.223)/20 = \qty{-0.116}{min^{-1}}$。

\textbf{10.} $k = \qty{0.116}{min^{-1}}$。

\textbf{11.} $0.693 / 0.116 = \qty{6.0}{min}$，与表中读数一致。

\textbf{12.} $k A_0 = 0.116 \times 0.800 = \qty{0.093}{min^{-1}}$，以
\omterm{def:g11:absorbance:absorbance}{吸光度}单位每分钟计。

\textbf{13.} 同样为 \qty{6.0}{min}：\omterm{def:g12:reaction-rates:first-order}{一级反应}的\omterm{def:g12:reaction-rates:half-life}{半衰期}
与初始值无关。

\textbf{14.} $0.800\, e^{-0.116 \times 30} = 0.025$。

\textbf{15.} $t = \ln 100 / k = 4.61 / 0.116 = \qty{40}{min}$。

\textbf{16.} $40 / 6.0 = 6.6$ 个\omterm{def:g12:reaction-rates:half-life}{半衰期}。

\textbf{17.} $t = \ln(0.800/0.010)/k = 4.38 / 0.116 = \qty{38}{min}$。

\textbf{18.} 仍为 6.6 个\omterm{def:g12:reaction-rates:half-life}{半衰期}，此时为 $6.6 \times 3.0 = \qty{20}{min}$。

\textbf{19.} 约 6.6 个\omterm{def:g12:reaction-rates:half-life}{半衰期}，即此处约 \qty{40}{min}。
\end{solution}
